% ============================================================
%  styles/preamble.tex
%  Shared preamble for Holonomic Space
% ============================================================

% ---- encoding and fonts -------------------------------------
\usepackage[T1]{fontenc}
\usepackage[utf8]{inputenc}
\usepackage{lmodern}
\usepackage{microtype}

% ---- mathematics --------------------------------------------
\usepackage{amsmath}
\usepackage{amssymb}
\usepackage{amsthm}
\usepackage{mathtools}
\usepackage{bm}           % bold math
\usepackage{bbm}          % blackboard bold

% ---- diagrams and commutative diagrams ----------------------
\usepackage{tikz}
\usepackage{tikz-cd}
\usetikzlibrary{arrows,positioning,calc}

% ---- layout -------------------------------------------------
\usepackage[
  top=1in, bottom=1in,
  inner=1.25in, outer=1in,
  headheight=14pt
]{geometry}
\usepackage{fancyhdr}
\pagestyle{fancy}
\fancyhf{}
\fancyhead[LE]{\small\leftmark}
\fancyhead[RO]{\small\rightmark}
\fancyfoot[C]{\thepage}

% ---- hyperlinks (load near last) ----------------------------
\usepackage[
  colorlinks=true,
  linkcolor=black,
  citecolor=black,
  urlcolor=blue
]{hyperref}

% ---- bibliography -------------------------------------------
\usepackage[
  backend=biber,
  style=alphabetic,
  sorting=nyt
]{biblatex}
\addbibresource{back/references.bib}

% ---- graphics -----------------------------------------------
\usepackage{graphicx}
\usepackage{float}

% ---- tables -------------------------------------------------
\usepackage{booktabs}
\usepackage{array}

% ---- enumeration --------------------------------------------
\usepackage{enumitem}

% ============================================================
%  Theorem environments
% ============================================================

\theoremstyle{plain}
\newtheorem{theorem}{Theorem}[chapter]
\newtheorem{proposition}[theorem]{Proposition}
\newtheorem{lemma}[theorem]{Lemma}
\newtheorem{corollary}[theorem]{Corollary}
\newtheorem{conjecture}[theorem]{Conjecture}

\theoremstyle{definition}
\newtheorem{definition}[theorem]{Definition}
\newtheorem{example}[theorem]{Example}
\newtheorem{construction}[theorem]{Construction}

\theoremstyle{remark}
\newtheorem{remark}[theorem]{Remark}
\newtheorem{note}[theorem]{Note}

% ============================================================
%  Custom commands
% ============================================================

% Fields
\newcommand{\vv}{\mathbf{v}}
\newcommand{\Ss}{S}
\newcommand{\Ll}{L}
\newcommand{\Dd}{D}

% Spaces
\newcommand{\Monoturn}{\mathcal{M}}
\newcommand{\Ocularum}{\mathcal{O}}
\newcommand{\Adm}{\mathcal{A}}
\newcommand{\Glob}{\mathbf{Glob}}
\newcommand{\Loc}{\mathbf{Loc}}
\newcommand{\Sob}[1]{H^{#1}}
\newcommand{\SobN}[2]{H^{#1}(#2)}

% Free energy
\newcommand{\FE}{\mathcal{F}}

% Functors
\newcommand{\Fdecomp}{F}
\newcommand{\Grecon}{G}

% Reconstruction error
\newcommand{\RecErr}[1]{E(#1)}

% Kernel notation
\newcommand{\kernel}{\ker}

% Cohomology
\newcommand{\Hone}{H^1}
\newcommand{\Hzero}{H^0}

% Coarsening length
\newcommand{\coarseningL}{\ell(t)}

% RSVP relaxation rate
\newcommand{\Rrsvp}{\mathcal{R}_{\mathrm{RSVP}}}

% Effective Hubble
\newcommand{\Heff}{H_{\mathrm{eff}}}

% Variational derivative shorthand
\newcommand{\varder}[2]{\frac{\delta #1}{\delta #2}}

% Norm shorthand
\newcommand{\norm}[1]{\left\|#1\right\|}
\newcommand{\normHk}[2]{\left\|#1\right\|_{H^{#2}}}

% Part-whole notation
\newcommand{\mso}{\mathrm{MSO}}

% Boxed principle (used for preface theorems and key results)
\newcommand{\principle}[1]{%
  \begin{center}
  \fbox{\parbox{0.85\textwidth}{\centering #1}}
  \end{center}
}

% Chapter epigraph (optional)
\newcommand{\epigraph}[2]{%
  \begin{flushright}
  \itshape #1\\[2pt]
  \upshape\small --- #2
  \end{flushright}
  \vspace{1em}
}

% Status labels for theorems
\newcommand{\statusproven}{\textbf{[Proven]}}
\newcommand{\statusconjectured}{\textbf{[Conjectured]}}
\newcommand{\statusopen}{\textbf{[Open Problem]}}

% ============================================================
%  Navigation box system
%  Requires tcolorbox
% ============================================================
\usepackage{tcolorbox}
\tcbuselibrary{breakable,skins}

% Colour palette
\definecolor{boxblue}{RGB}{220,232,245}
\definecolor{boxgreen}{RGB}{220,242,220}
\definecolor{boxorange}{RGB}{255,237,210}
\definecolor{boxred}{RGB}{250,220,220}
\definecolor{boxgray}{RGB}{235,235,235}
\definecolor{boxpurple}{RGB}{235,220,245}
\definecolor{boxyellow}{RGB}{255,248,210}
\definecolor{boxteal}{RGB}{210,240,238}

% ------ Why This Matters ------------------------------------
\newenvironment{WhyMatters}[1][]{%
  \begin{tcolorbox}[
    breakable,
    title={\textbf{Why This Matters}\ifx&#1&\else: #1\fi},
    colback=boxblue,colframe=blue!40!black,
    fonttitle=\small\sffamily,
    left=4pt,right=4pt,top=4pt,bottom=4pt
  ]
}{\end{tcolorbox}}

% ------ Technical Note --------------------------------------
\newcounter{technote}[chapter]
\renewcommand{\thetechnote}{\thechapter.\arabic{technote}}
\newenvironment{TechNote}[1][]{%
  \refstepcounter{technote}%
  \begin{tcolorbox}[
    breakable,
    title={\textbf{Technical Note~\thetechnote}\ifx&#1&\else: #1\fi},
    colback=boxgray,colframe=gray!60!black,
    fonttitle=\small\sffamily,
    left=4pt,right=4pt,top=4pt,bottom=4pt
  ]
}{\end{tcolorbox}}

% ------ Reconstruction Exercise -----------------------------
\newcounter{rcexercise}[chapter]
\renewcommand{\thercexercise}{\thechapter.\arabic{rcexercise}}
\newenvironment{ReconExercise}[1][]{%
  \refstepcounter{rcexercise}%
  \begin{tcolorbox}[
    breakable,
    title={\textbf{Reconstruction Exercise~\thercexercise}%
           \ifx&#1&\else: #1\fi},
    colback=boxgreen,colframe=green!50!black,
    fonttitle=\small\sffamily,
    left=4pt,right=4pt,top=4pt,bottom=4pt
  ]
}{\end{tcolorbox}}

% ------ Forward Reference -----------------------------------
\newenvironment{ForwardRef}[1][]{%
  \begin{tcolorbox}[
    breakable,
    title={\textbf{Forward Reference}\ifx&#1&\else: #1\fi},
    colback=boxyellow,colframe=orange!70!black,
    fonttitle=\small\sffamily,
    left=4pt,right=4pt,top=4pt,bottom=4pt
  ]
}{\end{tcolorbox}}

% ------ Status Box ------------------------------------------
\newenvironment{StatusBox}[1]{%
  \begin{tcolorbox}[
    breakable,
    title={\textbf{Status:} #1},
    colback=boxorange,colframe=orange!60!black,
    fonttitle=\small\sffamily,
    left=4pt,right=4pt,top=4pt,bottom=4pt
  ]
}{\end{tcolorbox}}

% ------ Common Misconception --------------------------------
\newenvironment{Misconception}[1][]{%
  \begin{tcolorbox}[
    breakable,
    title={\textbf{Common Misconception}\ifx&#1&\else: #1\fi},
    colback=boxred,colframe=red!50!black,
    fonttitle=\small\sffamily,
    left=4pt,right=4pt,top=4pt,bottom=4pt
  ]
}{\end{tcolorbox}}

% ------ Admissibility Preview -------------------------------
\newenvironment{AdmPreview}[1][]{%
  \begin{tcolorbox}[
    breakable,
    title={\textbf{Admissibility Preview}\ifx&#1&\else: #1\fi},
    colback=boxpurple,colframe=purple!50!black,
    fonttitle=\small\sffamily,
    left=4pt,right=4pt,top=4pt,bottom=4pt
  ]
}{\end{tcolorbox}}

% ------ Mathematical Echo -----------------------------------
\newenvironment{MathEcho}{%
  \begin{tcolorbox}[
    breakable,
    title={\textbf{Mathematical Echo}},
    colback=boxgray,colframe=gray!50!black,
    fonttitle=\small\sffamily,
    left=4pt,right=4pt,top=4pt,bottom=4pt
  ]
}{\end{tcolorbox}}

% ------ Open Question ---------------------------------------
\newcounter{openq}[chapter]
\renewcommand{\theopenq}{\thechapter.\arabic{openq}}
\newenvironment{OpenQuestion}[1][]{%
  \refstepcounter{openq}%
  \begin{tcolorbox}[
    breakable,
    title={\textbf{Open Question~\theopenq}\ifx&#1&\else: #1\fi},
    colback=boxteal,colframe=teal!60!black,
    fonttitle=\small\sffamily,
    left=4pt,right=4pt,top=4pt,bottom=4pt
  ]
}{\end{tcolorbox}}

% ------ Derivation Roadmap ----------------------------------
\newenvironment{DerivRoadmap}[1][]{%
  \begin{tcolorbox}[
    breakable,
    title={\textbf{Derivation Roadmap}\ifx&#1&\else: #1\fi},
    colback=boxyellow,colframe=yellow!60!black,
    fonttitle=\small\sffamily,
    left=4pt,right=4pt,top=4pt,bottom=4pt
  ]
}{\end{tcolorbox}}

% ------ Obstruction Box -------------------------------------
\newenvironment{ObstructionBox}[1][]{%
  \begin{tcolorbox}[
    breakable,
    title={\textbf{Obstruction}\ifx&#1&\else: #1\fi},
    colback=boxred,colframe=red!60!black,
    fonttitle=\small\sffamily,
    left=4pt,right=4pt,top=4pt,bottom=4pt
  ]
}{\end{tcolorbox}}

% ------ Interpretation Box ----------------------------------
\newenvironment{InterpBox}[1][]{%
  \begin{tcolorbox}[
    breakable,
    title={\textbf{Interpretation}\ifx&#1&\else: #1\fi},
    colback=boxblue,colframe=blue!50!black,
    fonttitle=\small\sffamily,
    left=4pt,right=4pt,top=4pt,bottom=4pt
  ]
}{\end{tcolorbox}}

% ------ Local-to-Global Insight -----------------------------
\newenvironment{LocalGlobal}[1][]{%
  \begin{tcolorbox}[
    breakable,
    title={\textbf{Local-to-Global Insight}\ifx&#1&\else: #1\fi},
    colback=boxgreen,colframe=green!60!black,
    fonttitle=\small\sffamily,
    left=4pt,right=4pt,top=4pt,bottom=4pt
  ]
}{\end{tcolorbox}}

% ------ Holonomy Note (end of each Part) --------------------
\newenvironment{HolonomyNote}{%
  \begin{tcolorbox}[
    breakable,
    title={\textbf{Holonomy Note}},
    colback=boxpurple,colframe=purple!60!black,
    fonttitle=\small\sffamily,
    left=4pt,right=4pt,top=4pt,bottom=4pt
  ]
}{\end{tcolorbox}}

% ------ Observable Consequences (Part VII) ------------------
\newenvironment{ObsConseq}[1][]{%
  \begin{tcolorbox}[
    breakable,
    title={\textbf{Observable Consequences}\ifx&#1&\else: #1\fi},
    colback=boxteal,colframe=teal!50!black,
    fonttitle=\small\sffamily,
    left=4pt,right=4pt,top=4pt,bottom=4pt
  ]
}{\end{tcolorbox}}

% ------ What Survived the Projection? -----------------------
%  Used as final section of every chapter.
%  Arg 1: Projection name
%  Arg 2: Preserved (tabular rows: Item & description \\)
%  Arg 3: Lost (tabular rows)
%  Arg 4: Reconstruction status
%  Arg 5: Kernel description
\newcommand{\WhatSurvived}[5]{%
  \section*{What Survived the Projection?}
  \addcontentsline{toc}{section}{What Survived the Projection?}
  \begin{tcolorbox}[
    breakable,
    title={\textbf{What Survived the Projection?}
           \quad\textit{Projection:} #1},
    colback=boxgray,colframe=gray!60!black,
    fonttitle=\small\sffamily
  ]
  \begin{minipage}[t]{0.48\linewidth}
    \textbf{Preserved}\\[2pt]
    \begin{tabular}{@{}l}
      #2
    \end{tabular}
  \end{minipage}
  \hfill
  \begin{minipage}[t]{0.48\linewidth}
    \textbf{Lost}\\[2pt]
    \begin{tabular}{@{}l}
      #3
    \end{tabular}
  \end{minipage}

  \medskip
  \textbf{Reconstruction possible:} #4

  \textbf{Kernel:} #5
  \end{tcolorbox}
}

% ------ Cross-Domain Example table --------------------------
\newenvironment{CrossDomain}[1][]{%
  \begin{tcolorbox}[
    breakable,
    title={\textbf{Same Mathematics, Different Systems}%
           \ifx&#1&\else: #1\fi},
    colback=boxblue,colframe=blue!40!black,
    fonttitle=\small\sffamily
  ]
}{\end{tcolorbox}}

% ============================================================
%  Extended navigation boxes (added in v6)
% ============================================================

% ------ Why This Chapter Exists ----------------------------
\newenvironment{WhyChapter}[1][]{%
  \begin{tcolorbox}[
    breakable,
    title={\textbf{Why This Chapter Exists}%
           \ifx&#1&\else: #1\fi},
    colback=boxblue,colframe=blue!50!black,
    fonttitle=\small\sffamily,
    left=4pt,right=4pt,top=4pt,bottom=4pt
  ]
}{\end{tcolorbox}}

% ------ What Has Been Assumed -------------------------------
\newenvironment{Assumptions}[1][]{%
  \begin{tcolorbox}[
    breakable,
    title={\textbf{Assumptions Behind This Derivation}%
           \ifx&#1&\else: #1\fi},
    colback=boxorange,colframe=orange!50!black,
    fonttitle=\small\sffamily,
    left=4pt,right=4pt,top=4pt,bottom=4pt
  ]
}{\end{tcolorbox}}

% ------ Connection to Earlier Chapters ----------------------
\newenvironment{ChapterConnection}[1][]{%
  \begin{tcolorbox}[
    breakable,
    title={\textbf{Connection to Earlier Chapters}%
           \ifx&#1&\else: #1\fi},
    colback=boxteal,colframe=teal!50!black,
    fonttitle=\small\sffamily,
    left=4pt,right=4pt,top=4pt,bottom=4pt
  ]
}{\end{tcolorbox}}

% ------ Not Unique to Cosmology ----------------------------
\newenvironment{NotUnique}[1][]{%
  \begin{tcolorbox}[
    breakable,
    title={\textbf{Why This Is Not Unique to Cosmology}%
           \ifx&#1&\else: #1\fi},
    colback=boxgreen,colframe=green!50!black,
    fonttitle=\small\sffamily,
    left=4pt,right=4pt,top=4pt,bottom=4pt
  ]
}{\end{tcolorbox}}

% ------ What Could Falsify This? ---------------------------
\newenvironment{Falsifiability}[1][]{%
  \begin{tcolorbox}[
    breakable,
    title={\textbf{What Could Falsify This?}%
           \ifx&#1&\else: #1\fi},
    colback=boxred,colframe=red!50!black,
    fonttitle=\small\sffamily,
    left=4pt,right=4pt,top=4pt,bottom=4pt
  ]
}{\end{tcolorbox}}

% ------ Reconstruction View --------------------------------
\newenvironment{ReconView}[1][]{%
  \begin{tcolorbox}[
    breakable,
    title={\textbf{Reconstruction View}%
           \ifx&#1&\else: #1\fi},
    colback=boxpurple,colframe=purple!50!black,
    fonttitle=\small\sffamily,
    left=4pt,right=4pt,top=4pt,bottom=4pt
  ]
}{\end{tcolorbox}}

% ------ Worked Example -------------------------------------
\newcounter{workedex}[chapter]
\renewcommand{\theworkedex}{\thechapter.\arabic{workedex}}
\newenvironment{WorkedExample}[1][]{%
  \refstepcounter{workedex}%
  \begin{tcolorbox}[
    breakable,
    title={\textbf{Worked Example~\theworkedex}%
           \ifx&#1&\else: #1\fi},
    colback=boxyellow,colframe=yellow!60!black,
    fonttitle=\small\sffamily,
    left=4pt,right=4pt,top=4pt,bottom=4pt
  ]
}{\end{tcolorbox}}

% ------ Problem Set (end of chapter) -----------------------
%  Usage: \begin{ProblemSet} \problem{...} \problem{...} \end{ProblemSet}
\newcounter{problem}[chapter]
\renewcommand{\theproblem}{\thechapter.\arabic{problem}}
\newcommand{\problem}[1]{%
  \refstepcounter{problem}%
  \medskip\noindent\textbf{Problem~\theproblem.}\quad #1\par
}
\newenvironment{ProblemSet}{%
  \section*{Problems}
  \addcontentsline{toc}{section}{Problems}
}{\medskip}

% ------ Visual Note (diagram reference) --------------------
\newenvironment{VisualNote}[1][]{%
  \begin{tcolorbox}[
    breakable,
    title={\textbf{Visual Guide}%
           \ifx&#1&\else: #1\fi},
    colback=boxgray,colframe=gray!40!black,
    fonttitle=\small\sffamily,
    left=4pt,right=4pt,top=4pt,bottom=4pt
  ]
}{\end{tcolorbox}}

% ------ Observable Box (Chapter 63 final form) -------------
\newenvironment{ObservableBox}[1][]{%
  \begin{tcolorbox}[
    breakable,
    title={\textbf{Observable vs.\ Reconstructible}%
           \ifx&#1&\else: #1\fi},
    colback=boxblue,colframe=blue!60!black,
    fonttitle=\small\sffamily,
    left=4pt,right=4pt,top=4pt,bottom=4pt
  ]
}{\end{tcolorbox}}

% ============================================================
%  Chapter opening prose interlude
%  Appears after epigraph, before first section
% ============================================================
\newenvironment{ChapterIntro}{%
  \medskip
  \begin{quote}\itshape
}{\end{quote}\medskip}

% ============================================================
%  Further Reading box and citation integration
% ============================================================
\newenvironment{FurtherReading}{%
  \begin{tcolorbox}[
    breakable,
    title={\textbf{Further Reading}},
    colback=boxgray,colframe=gray!50!black,
    fonttitle=\small\sffamily,
    left=4pt,right=4pt,top=4pt,bottom=4pt
  ]
}{\end{tcolorbox}}

% Citation anchor note (appears after a theorem to place it in literature)
\newenvironment{CitationAnchor}{%
  \begin{tcolorbox}[
    breakable,
    title={\textbf{Mathematical Ancestry}},
    colback=boxyellow,colframe=yellow!50!black,
    fonttitle=\small\sffamily,
    left=4pt,right=4pt,top=4pt,bottom=4pt
  ]
}{\end{tcolorbox}}

% usetikzlibrary additions needed for diagrams
\usetikzlibrary{decorations.pathmorphing,patterns}
